% ============================================================================
% Section 6: Non-Blocking FSM (Task 6)
% Handout s10.2 requires ONLY: final state-machine diagram, short description
% of the transition logic, why the implementation is non-blocking, one example
% write/read/verify result.
% The diagram must show (Task 6 "State-machine record"): initial state,
% conditions causing transitions, success path, failure path, abort path.
% State numbers match ee_state_t in task6_fsm.c.
% ============================================================================
\section{Non-Blocking FSM}
\label{sec:fsm}

\begin{figure}[htbp]
  \centering
  \resizebox{\linewidth}{!}{%
  \begin{tikzpicture}[
      >={Stealth[length=2mm]},
      st/.style={draw, rounded corners, minimum height=10mm, minimum width=17mm,
                 align=center, font=\footnotesize},
      lbl/.style={font=\scriptsize, align=center},
      jn/.style={circle, fill, inner sep=0.9pt}]
    % ---- states on a grid (3.5 cm between centres) -------------------------
    \node[st, initial, initial text={reset}] (idle) at (0,0)    {IDLE\\(0)};
    \node[st]                     (prep)   at (3.5,0)  {WRITE\\PREP (1)};
    \node[st]                     (write)  at (7.0,0)  {WRITE\\(2)};
    \node[st]                     (busy)   at (10.5,0)  {WAIT\\BUSY (3)};
    \node[st]                     (read)   at (14.0,0) {READ\\(4)};
    \node[st]                     (verify) at (17.5,0) {VERIFY\\(5)};
    \node[st, fill=green!12]      (ok)     at (21.0,0) {SUCCESS\\(6)};
    \node[st, fill=red!12]        (fail)   at (21.0,2.8) {FAIL\\(7)};

    % ---- main (success) path ------------------------------------------------
    \draw[->] (idle)   -- node[lbl, above] {PA0} (prep);
    \draw[->] (prep)   -- node[lbl, above] {RDY=0\\WREN} (write);
    \draw[->] (write)  -- node[lbl, above] {WEL=1\\WRITE} (busy);
    \draw[->] (busy)   -- node[lbl, above] {RDY=0} (read);
    \draw[->] (read)   -- node[lbl, above] {byte\\read} (verify);
    \draw[->] (verify) -- node[lbl, above, pos=0.7] {match} (ok);

    % ---- WAIT_BUSY self-loop (square) ---------------------------------------
    \draw[->] ([xshift=-4mm]busy.south) -- ++(0,-6mm) -| ([xshift=4mm]busy.south);
    \node[lbl, below=6.5mm of busy.south] (looplbl) {RDY=1: RDSR\\every 1\,ms};

    % ---- failure paths: up to a common line into FAIL -----------------------
    \draw[->] (prep.north |- fail.west) -- (fail.west);
    \draw (prep.north)   -- node[lbl, right, pos=0.55] {bad $A$ or\\busy timeout} (prep.north |- fail.west);
    \draw (write.north)  -- node[lbl, right, pos=0.55] {WEL=0} (write.north |- fail.west);
    \draw (busy.north)   -- node[lbl, right, pos=0.55] {timeout\\50\,ms} (busy.north |- fail.west);
    \draw (verify.north) -- node[lbl, right, pos=0.55] {mismatch} (verify.north |- fail.west);
    \node[jn] at (write.north |- fail.west) {};
    \node[jn] at (busy.north |- fail.west) {};
    \node[jn] at (verify.north |- fail.west) {};

    % ---- PA0 restart from SUCCESS / FAIL: along the bottom into WRITE_PREP --
    \coordinate (rail) at (0,-3.4);
    \draw (fail.east) -- ++(6mm,0) coordinate (fe) -- (fe |- rail);
    \draw (ok.south) -- (ok.south |- rail);
    \node[jn] at (ok.south |- rail) {};
    \draw[->] (fe |- rail) -- node[lbl, below, pos=0.5] {PA0 (new write)} (prep.south |- rail) -- (prep.south);

    % ---- abort: dashed box of busy states -> IDLE ---------------------------
    \begin{scope}[on background layer]
      \node[draw, dashed, rounded corners, inner xsep=4mm, inner ysep=8mm,
            fit=(prep)(verify)(looplbl)] (box) {};
    \end{scope}
    \draw[->, dashed] ([xshift=3mm]box.south west) -- ++(0,-6mm) coordinate (ab)
      -- node[lbl, below] {PA3 (abort)} (idle.south |- ab) -- (idle.south);
  \end{tikzpicture}}
  \caption{EEPROM transaction state machine (\reg{ee\_state} values in
    brackets). Dashed box: busy states, all aborted by PA3.}
  \label{fig:fsm}
\end{figure}

\paragraph{Transition logic.}
PA0 starts a transaction from IDLE, SUCCESS or FAIL; a press while busy is
ignored. WRITE\_PREP waits for RDY~=~0 and sends WREN; WRITE checks WEL~=~1 and
sends WRITE, two address bytes and $B$; WAIT\_BUSY reads the status once per
\SI{1}{\milli\second} until RDY~=~0~\cite[p.~5, Tab.~10]{eeprom}; READ reads
$A$ back and VERIFY compares it with $B$. A busy timeout, WEL~=~0 or a mismatch
leads to FAIL (red), a match to SUCCESS (green). PA3 returns any busy state to
IDLE with both status LEDs off.

\paragraph{Why it is non-blocking.}
Each call of \reg{update\_eeprom\_state\_machine()} performs at most one SPI
frame (at most \SI{128}{\micro\second} of clocks) and returns. While the
EEPROM is writing, the FSM only compares \reg{HAL\_GetTick()} with the time of
the last status check and returns if the next check is not due, so the main
loop keeps running \reg{read\_inputs()} and PA3 is handled on the next pass.
Every SPI frame finishes inside its own call, so an abort always leaves CS high
and the bus idle.

\paragraph{Example result.}
$B = \text{\BHex}$ written to $A = \text{\AHex}$: the write finished after
4 status checks (\SI{4}{\milli\second}) while the main loop ran 112 times; read
back as 0x44; verification passed.
